\documentclass{article}

\usepackage[preprint]{tmlr}
% \usepackage{tmlr}            % anonymized submission
% \usepackage[accepted]{tmlr}  % camera-ready

\usepackage{amsmath,amssymb,amsfonts}
\usepackage{booktabs}
\usepackage{graphicx}
\usepackage{microtype}
\usepackage{hyperref}
\usepackage{url}

\title{Auditable Compositional Question Answering over UK Public Procurement}

% TODO final author block
\author{\name Anonymous \email anonymous@example.com \\
 \addr University College London}

\def\month{07}
\def\year{2026}
\def\openreview{\url{https://openreview.net/forum?id=XXXX}}

\begin{document}

\maketitle

\begin{abstract}
The United Kingdom publishes hundreds of thousands of contract award
notices, yet the public cannot reliably analyse them in natural language,
because the questions that matter are exhaustive aggregations where a
fluent wrong answer is worse than none. We introduce, to the best of our
knowledge, the first systematic benchmark and executable framework for
auditable compositional question answering over UK public procurement. The
task ships complete. A convention-explicit knowledge graph of 215{,}221
contract awards, a benchmark of 12{,}828 questions each carrying an
executable gold program, a typed compositional operator algebra that
formalises real procurement analytics, refusal on ambiguous, unsupported
and empty questions as scored behaviour, and answer oracles validated by
an independent second implementation at 99.88 percent over 14{,}770
audited answers. Because no annotator pool can label compositional
procurement programs at scale, training uses a verified bootstrap. Models
propose candidate programs and deterministic checks, compilation,
grounding, execution and oracle agreement, decide what may be learned
from. An 8-billion-parameter local model trained this way reaches 85.65
percent against 69.76 for the cloud teacher that generated its training
data (paired McNemar $p < 10^{-15}$), and 78.31 percent on a sealed
challenge set with an independent surface grammar, a constraint-necessity
audit, and a 28-point margin over the same teacher. Transplanted to
WikiTableQuestions by exchanging only the data-representation layer, the
recipe lifts a 22.5 percent zero-shot floor to 44.4 with answer-only
supervision and 51.8 with gold programs on the official sealed test,
while a four-way attribution audit shows a substantial share of the
apparent expressiveness gap to be representation and compilation debt
rather than missing reasoning capability. We claim no state of the art.
The claim is a task built end to end, and audits strong enough to say
precisely what its system can and cannot yet do.
\end{abstract}

% ---------------------------------------------------------------
\section{Introduction}\label{sec:intro}
% TODO convert docs/final_paper/draft_01_introduction.md

\section{Related Work}\label{sec:related}
% TODO convert docs/final_paper/draft_02_related.md

\section{The Task and the System}\label{sec:system}
% TODO convert docs/final_paper/draft_03_system.md
% Figure 1 (architecture) goes here.

\section{Benchmarks and Audits}\label{sec:benchmarks}
% TODO convert docs/final_paper/draft_04_benchmarks.md

\section{Experiments}\label{sec:experiments}
% TODO convert docs/final_paper/draft_05_experiments.md
% Figure 2 (scoreboard ladder), Figure 3 (attribution matrix), Tables 1-3.

\section{Discussion}\label{sec:discussion}
% TODO convert docs/final_paper/draft_06_07_discussion_conclusion.md (§6)

\section{Conclusion}\label{sec:conclusion}
% TODO convert docs/final_paper/draft_06_07_discussion_conclusion.md (§7)

\bibliography{references}
\bibliographystyle{tmlr}

\appendix
\section{Operator Algebra Reference}\label{app:algebra}
% TODO
\section{Audit Protocols}\label{app:audits}
% TODO
\section{Prompt and Serving Configuration}\label{app:config}
% TODO

\end{document}
